\section*{الفصل \ref{ch:g9:elements-universe-earth} --- العناصر في الكون وعلى الأرض}

\begin{solution}{exo:g9:elements-universe-earth:1}
الهيدروجين (نحو \qty{74}{\percent}) والهيليوم (نحو
\qty{25}{\percent}).
\end{solution}

\begin{solution}{exo:g9:elements-universe-earth:2}
الأكسجين في الحالتين.
\end{solution}

\begin{solution}{exo:g9:elements-universe-earth:3}
داخل النجوم، وفي انفجاراتها.
\end{solution}

\begin{solution}{exo:g9:elements-universe-earth:4}
نحو \qty{3.5}{\percent}.
\end{solution}

\begin{solution}{exo:g9:elements-universe-earth:5}
$46.6 + 27.7 = \qty{74.3}{\percent}$: أي نحو ثلاثة أرباع.
\end{solution}

\begin{solution}{exo:g9:elements-universe-earth:6}
فهما غازان خفيفان جدًا: أفلتا من الأرض في شبابها إلى الفضاء (والهيليوم لا
يكوّن مركبات كان يمكن أن تحبسه في الصخور).
\end{solution}

\begin{solution}{exo:g9:elements-universe-earth:7}
$0.23 \times 70 = \qty{16.1}{kg}$.
\end{solution}

\begin{solution}{exo:g9:elements-universe-earth:8}
\omterm{def:g7:atoms-and-molecules:atom}{ذرات} الأكسجين أثقل من \omterm{def:g7:atoms-and-molecules:atom}{ذرات} الهيدروجين بست عشرة مرة: فعدد أقل من \omterm{def:g7:atoms-and-molecules:atom}{ذرات}
الأكسجين يزن مع ذلك أكثر.
\end{solution}

\begin{solution}{exo:g9:elements-universe-earth:9}
نحو $0.081 \times 1000 = \qty{81}{kg}$ من الألومنيوم.
\end{solution}

\begin{solution}{exo:g9:elements-universe-earth:10}
الأكسجين والكالسيوم (وخارج العناصر الثمانية الرئيسية للقشرة، يوجد الهيدروجين
والكربون والفوسفور في القشرة بمقادير
صغيرة).
\end{solution}

\begin{solution}{exo:g9:elements-universe-earth:11}
$40 \times 1.67 \times 10^{-27} = \qty{6.68e-26}{kg}$؛
و$1.06 / 6.68 \times 10^{-26} \approx 1.6 \times 10^{25}$ \omterm{def:g7:atoms-and-molecules:atom}{ذرة} كالسيوم.
\end{solution}

\begin{solution}{exo:g9:elements-universe-earth:12}
\omterm{def:g7:atoms-and-molecules:atom}{الذرات} لا تُخلق ولا تفنى بالتغيرات الكيميائية: بل تنتقل فقط من مركب إلى آخر
(صخر، وغاز، وسكر، وجسم حي). \omterm{def:g7:atoms-and-molecules:atom}{فذرات} الكربون على الأرض تُستعمل مرة بعد
مرة.
\end{solution}

\begin{solution}{pb:g9:elements-universe-earth:1}
\textbf{1.} $0.61 \times 70 = \qty{42.7}{kg}$.

\textbf{2.} الكربون $0.23 \times 70 = \qty{16.1}{kg}$؛ والهيدروجين
$0.10 \times 70 = \qty{7.0}{kg}$.

\textbf{3.} $61 + 23 + 10 = \qty{94}{\percent}$.

\textbf{4.} الماء، \ce{H2O}.

\textbf{5.} $16 \times 1.67 \times 10^{-27} = \qty{2.67e-26}{kg}$.

\textbf{6.} الكربون: $12 \times 1.67 \times 10^{-27} = \qty{2.00e-26}{kg}$؛
والهيدروجين: \qty{1.67e-27}{kg}.

\textbf{7.} 16.

\textbf{8.} \omterm{def:g9:inside-the-atom:nucleus}{الإلكترون} أخف من \omterm{def:g9:inside-the-atom:proton}{النوكليون} بنحو 1836 مرة: \omterm{def:g9:inside-the-atom:nucleus}{فالإلكترونات} أقل من جزء
من ألف من الكتلة.

\textbf{9.} $7.0 / 1.67 \times 10^{-27} \approx 4.2 \times 10^{27}$
\omterm{def:g7:atoms-and-molecules:atom}{ذرة} هيدروجين.

\textbf{10.} $16.1 / 2.00 \times 10^{-26} \approx 8.0 \times 10^{26}$
\omterm{def:g7:atoms-and-molecules:atom}{ذرة} كربون.

\textbf{11.} $42.7 / 2.67 \times 10^{-26} \approx 1.60 \times 10^{27}$
\omterm{def:g7:atoms-and-molecules:atom}{ذرة} أكسجين.

\textbf{12.} حسابنا، $1.60 \times 10^{27}$، يتفق مع التقدير
$1.61 \times 10^{27}$: فجسم كتلته \qty{70}{kg} يحوي نحو $1.6 \times
10^{27}$ \omterm{def:g7:atoms-and-molecules:atom}{ذرة} أكسجين.
\end{solution}
